\uuid{GXfJ}
\exo7id{4230}
\titre{exo7 4230}
\auteur{quercia}
\organisation{exo7}
\datecreate{2010-03-12}
\isIndication{false}
\isCorrection{false}
\chapitre{Calcul d'intégrales}
\sousChapitre{Somme de Riemann}
\cpge{mpsi-mp2i.nombres-reels-suites.d, pcsi.nombres-reels-suites.c, ptsi.nombres-reels-suites.c, mpsi-mp2i.limites-continuite.b, pcsi.limites-continuite.b, ptsi.limites-continuite.b, mpsi-mp2i.integration.c, pcsi.integration.b, ptsi.integration.b, bcpst1.suites-reelles, bcpst1.limites-continuite.c, bcpst1.integration-segment.a}
\module{Analyse}
\niveau{L1}
\difficulte{}

\contenu{
\question{
Soit $f : {[0,1]} \to \R$ continue.
Montrer que $\Bigl(1+\frac1nf\bigl(\frac1n\bigr)\Bigr)
     \Bigl(1+\frac1nf\bigl(\frac2n\bigr)\Bigr)\dots
     \Bigl(1+\frac1nf\bigl(\frac nn\bigr)\Bigr)
     \to \exp \int_{t=0}^1 f(t)\,d t$ lorsque $n\to\infty$.

(On pourra utiliser : $\forall\ x\ge -\frac12,\ x-x^2 \le \ln x \le x$)
}
}
